\documentclass[10pt]{article}
\usepackage[diary]{aditya}

\title{Diary entry for Spectra}
\author{Aditya Pahuja}
\date{\today}

\begin{document}
\maketitle
\tableofcontents
\pagebreak

\section{What are spectra?}
(Literal question, not philosophical.)
\begin{definition}[Spectrum]
    \label{def:spectrum}
    A \vocab{spectrum} $X$ is a sequence of based spaces $\{X_n\}_{n\ge0}$
    and a sequence of maps $\xi_n\colon\Sigma X_n\to X_{n+1}$,
    called \vocab{bonding maps}. A morphism of spectra
    $f\colon X\to Y$ is a sequence of maps $f_n\colon X_n\to Y_n$
    with the obvious commutativity property:
    \begin{center}
        \begin{tikzcd}
	    \Sigma X_n\arrow[r,"\xi_n"] \arrow[d,"\Sigma f_n"]
	    & X_{n+1}\arrow[d, "f_{n+1}"] \\
	    \Sigma Y_n\arrow[r, "\eta_n"] & Y_{n+1}
        \end{tikzcd}
    \end{center}
    (This definition is sometimes used instead for so-called ``prespectra''.
    If we need to distinguish these from other kinds of spectra
    we will call them ``sequential spectra''.)
\end{definition}
Note that we will always work with based spaces here
unless otherwise specified.

\begin{example}
    [Zero spectrum]
    \label{ex:zero-spectrum}
    The zero spectrum $*$ is the spectrum whose maps
    spaces are all the singleton spaces, which we will also denote as $*$.
    Since $\Sigma(*) = *$, there is exactly one choice for the bonding maps.
    The zero map between two spectra $X$ and $Y$
    is the unique map factoring through $*$,
    i.e. sending each $X_n$ to the basepoint of $Y_n$.
\end{example}

\begin{example}
    [Suspension spectra]
    \label{ex:suspension-spectrum}
    The sphere spectrum $\mathbb S$ is the spectrum
    whose $n$th space is the sphere $S^n$
    and whose bonding maps $\xi_n\colon \Sigma S^n\to S^{n+1}$
    are each the identity map, noting that
    $\Sigma S^n = S^n\wedge S^1 \cong S^{n+1}$.
    More generally, for a space $X$, the suspension spectrum
    $\Sigma^\infty X$ is the spectrum whose $n$th space is
    $\Sigma^n X$ and the bonding maps
    $\Sigma\Sigma^n X = \Sigma^{n+1}X\to \Sigma^{n+1}X$ are the identity maps.
\end{example}

\begin{definition}[Stable homotopy groups]
    \label{def:stable-homotopy-groups}
    For a spectrum $X$, the $k$th \vocab{stable homotopy group}
    of $X$, denoted $\pi_k(X)$ for an integer $k$,
    is the colimit of the system of abelian groups
    \begin{equation*}
	\cdots\too
	\pi_{n+k}(X_n)\overset\sigma\too \pi_{n+1+k}(\Sigma X_n)
	\overset{\xi_n}\too \pi_{n+1+k}(X_{n+1})\to\cdots
    \end{equation*}
    where $\sigma$ is the map sending $\varphi\colon S^{n+k}\to X_n$
    to the map $\sigma(\varphi)\colon S^{n+k+1}\to \Sigma X_n$
    defined by $\sigma(\varphi)(x,t) = \varphi(x)$.
    (If $k$ is negative then the system is defined starting at $n \ge -k$).
\end{definition}

To emphasize the last sentence: we can have
\emph{negative dimensional homotopy groups}!
In some sense, spectra let us think about ``negative dimensional geometry''.

\begin{example}
    [Shifted spectra]
    \label{ex:shifted-spectra}
    The colimit defining the $k$-th stable homotopy group
    does not change if one deletes finitely many terms (say the first $N$)
    from the directed system. Therefore, shifting a spectrum up or down
    (filling in with singleton spaces or truncating the spectrum
    respectively to account for indexing) just shifts the homotopy groups.
    In light of this, for a space $X$ and any integer $d$,
    define the spectrum $\Sigma^{\infty - d}X$
    to be the suspension spectrum shifted up $d$ spaces
    (so down $\abs d$ spaces if $d < 0$),
    i.e. $(\Sigma^{\infty - d}X)_n = (\Sigma X)_{n-d}$,
    where $(\Sigma X)_{k} = *$ if $k < 0$.
    The homotopy groups of $\Sigma^{\infty - d}A$ are
    $\pi_k(\Sigma^{\infty - d}A) = \pi_{k-d}(\Sigma^\infty A)$.

    The $d$-sphere $\mathbb S^d$ is then defined as
    $\Sigma^{\infty - d} S^0$, i.e. the spectrum
    which, when suspended $d$ times, is the usual sphere spectrum
    of \Cref{ex:suspension-spectrum}.
\end{example}

A \vocab{degree $k$ element} of a spectrum $X$
is a map $\varphi\colon S^{n+k}\to X_n$,
up to identification of $\varphi$ with $\sigma(\varphi)$,
\question{what does this identification mean}
and the stable homotopy groups are
the degree $k$ elements of $X$ up to homotopy.

\begin{definition}[Equivalences of spectra]
    \label{def:equivalences}
    A \vocab{stable equivalence} of spectra
    is a map of spectra inducing isomorphisms on every homotopy groups.
    Two spectra $X$ and $Y$ are \vocab{stably equivalent}
    if there is a sequence of spectra $X = X_1$, $X_2$, \dots, $X_{2n} = Y$
    and stable equivalences $X_1\to X_2$, $X_3\to X_2$, $X_3\to X_4$,
    \dots, $X_{2n-1}\to X_{2n}$ (the parity doesn't change anything
    since one can pad out the sequence with copies of $Y$).
    One says that $X$ and $Y$ are connected by
    a zig-zag of stable equivalences.

    A \vocab{level equivalences} is a map of spectra $X\to Y$
    in which each map $X_n\to Y_n$ is a weak equivalence.

    A \vocab{homotopy equivalence} is a map of spectra $f\colon X\to Y$
    such that there is a map $g\colon Y\to X$
    with the property that $g\circ f\simeq\operatorname{id}_X$
    and $f\circ g\simeq\operatorname{id}_Y$ (we will define homotopy later).
\end{definition}

We can now define stable homotopy theory:
\vocab{stable homotopy theory} is the study of
spectra up to stable equivalence.

\begin{example}
    [Inclusion as stable equivalence]
    \label{ex:inclusion-is-stable-equivalence}
    For a spectrum $X$, let $X'$ be the spectrum
    formed by replacing the first $d$ spaces with singleton sets.
    The inclusion $X'\into X$ (i.e. the map on each level is the inclusion)
    is a stable equivalence.
\end{example}

\begin{example}
    [Weak contractibility]
    \label{ex:weakly-contractible}
    For any spectrum $X$ there are unique maps
    $*\to X\to *$; if one of these (equivalently both)
    are stably equivalences, then $X$ is said to be weakly contractible.
    Note that this depends not only on the spaces in the spectrum
    but also on the bonding maps; for example,
    any spectrum can be made weakly contractible
    by setting the bonding maps to the maps sending $\Sigma X_n$
    to the basepoint of $X_{n+1}$.
\end{example}

\begin{example}
    [Thom spectra]
    \label{ex:thom-spectra}
    Recall that the \vocab{Thom space} $T(E)$
    of a (real) rank $n$ vector bundle $p\colon E\to B$
    is formed by compactifying each fiber to create an $n$-sphere bundle
    and then identifying all the points at infinity.
    Equivalently, $T(E)$ is homeomorphic to $D(E)/S(E)$,
    where, for a given metric on the vector bundle,
    $D(E)$ is the unit disk bundle and $S(E) = \partial D(E)$
    is the unit sphere bundle.
    For example, the Thom space of a trivial bundle
    $p\colon \RR^n\times B\to B$ is homeomorphic to
    $B_+\wedge S^n = \Sigma^n_+B$, where $B_+$ is the space
    $B\cup\{*\}$ with basepoint $*$
    (thus in general the Thom space is a ``twisted suspension'' of the base).

    Letting $\eps^q$ be the trivial rank $q$ vector bundle over $B$
    and $E$ any other vector bundle,
    observe that $T(E +\eps^q) = T(E\times\RR^q) = \Sigma^q T(E)$.
    \todo{prove vb fact}

    For another vector bundle $E'$, let $E''$ be such that
    $E' + E'' = \eps^q$, and form
    the formal difference $E - E' = E + E'' - E'' - E' = E + E'' - \eps^q$.
\end{example}


\end{document}

